\section{Phạm vi nghiên cứu}
Phạm vi của luận văn bao hàm việc khảo sát và hệ thống hóa các kiến thức nền tảng về giải tích lồi,
tối ưu hóa nón, lý thuyết đồ thị và bài toán đường đi ngắn nhất để làm cơ sở xây dựng GCS framework. 
Nghiên cứu tập trung sâu vào các phương pháp phân hoạch không gian tự do như Phân hoạch lồi xấp xỉ (ACD), 
kỹ thuật Iterative Regional Inflation by Semi-Definite Programming (IRIS) và 
Visibility Clique Cover (VCC) để xây dựng đồ thị tập các vùng an toàn. 
Trong khía cạnh động lực học, luận văn giới hạn việc mô hình hóa các ràng buộc vi phân liên tục bao gồm vận tốc, gia tốc và độ trơn (jerk) 
thông qua việc tận dụng tính chất bao lồi và các đạo hàm của đường cong Bézier. 
Các thực nghiệm số và đánh giá hiệu suất được thực hiện trên một loạt các môi trường có độ phức tạp tăng dần, từ các vật cản 2D và mê cung 
đến các hệ thống động lực học phức tạp như thiết bị bay quadrotor và tay máy 7 bậc tự do (7-DoF). 
Quá trình đánh giá được thực hiện dựa trên các chỉ số định lượng cụ thể như thời gian giải (solve time), chi phí quỹ đạo (path cost) và cận sai số tối ưu (optimality gap).